\chapter{Introduction} \label{chap:introduction}
%%
The demand for computational resources in modern scientific and engineering applications continues to grow, making parallel computing an essential area of development. While processor performance continues to improve, the growth in single-core performance has slowed considerably since the end of Dennard scaling~\cite{hennessy_new_2019}, making it impractical to tackle the largest problems on a single processor alone. Instead, it has become necessary to distribute scientific and engineering applications across the many processors of a parallel computing system.

To facilitate this, many models have been proposed that represent scientific and engineering applications as graphs or hypergraphs~\cite{umpa_2013, CatalyurekU.V.1999Hdfp, two_dimensional_part2010, hendrickson_graph_2000, WalshawC.1997PDGP}. The vertices correspond to units of computation and the edges encode data dependencies among them~\cite{hendrickson_graph_2000}. Efficiently distributing the computation then becomes a question of how to partition this graph among processors: balancing the workload evenly while minimizing inter-processor communication. This is also known as the \emph{graph partitioning problem}. Although graph partitioning is known to be NP-hard, a variety of algorithms have been proposed that produce good quality solutions in practice by using heuristics~\cite{karypis_analysis_1995, karypis_multilevelk-way_1998, karypis_parallel_1999, hendrickson_multilevel_1995}. Beyond parallel computing, graph partitioning plays a critical role in numerous other application areas, including task scheduling, neural network simulations, operations research and more.
%In 2005, Herb Sutter published \textit{Software and the Concurrency Revolution}~\cite{SutterHerb2005Satc}, arguing 

More than 20 years ago, the scientific community already argued that the future of software development lies in concurrency, e.g.,~\cite{SutterHerb2005Satc}. As computational demands continue to grow, our programs and requirements consistently outpace hardware improvements. Sutter and Larus highlight the need for better parallel programming tools, given the inherent difficulty of developing parallel solutions~\cite{SutterHerb2005Satc}. A decade later, on July 29, 2015, President Obama issued Executive Order 13702, establishing the National Strategic Computing Initiative (NSCI) to ensure continued United States leadership in high-performance computing~\cite{obama_nsci_2015}, underscoring the importance of this field and the need for sustained research. Graph partitioning algorithms are central to making this possible.

Most modern graph partitioning algorithms minimize \emph{edge-cut} (see Section~\ref{subs:objective_functions}), seeking to minimize the number of inter-partition edges. The appeal of this metric lies in its simplicity: it is easy to compute, and for the structured meshes that dominate scientific computing, it correlates reasonably well with actual communication costs~\cite{hendrickson_graph_2000}. However, edge-cut has several well-known shortcomings as a representation for communication~\cite{hendrickson_graph_2000, HENDRICKSONB1998Gpap}. First, it over-counts \textit{communication volume} (see Section~\ref{subs:objective_functions}): a vertex with multiple edges crossing into the same partition needs to be communicated only once, yet each edge is counted separately when using edge-cut. Second, and most relevant to this work, minimizing the \emph{total} edge-cut does not account for how communication is distributed across partitions. In a parallel application, performance is generally limited by its slowest process~\cite{hendrickson_graph_2000}, meaning that even if the total edge-cut is good, some process might bear a greater communication load than others, creating bottlenecks. 

Karypis and Moulitsas highlighted the need for partitioning algorithms that account for heterogeneous architectures and proposed using \emph{communication volume} as an optimization objective function instead of edge-cut~\cite{moulitsas_architecture_2008}. This communication volume objective function was subsequently added into the METIS software package (see Section~\ref{sec:Metis very general description}). However, like the edge-cut objective function, METIS focuses on the total communication volume without considering how volume is distributed among partitions. Regardless of the hardware architecture, focusing only on the total volume of a partitioned graph may lead to bottlenecks in some partitions.

Efficiently distributing computation across processors is central to high-performance computing, and graph partitioning is the standard technique for achieving this. In \textit{Graph Partitioning and Parallel Solvers: Has the Emperor No Clothes?}~\cite{HENDRICKSONB1998Gpap}, Hendrickson concludes that one of the key needs is to develop a more accurate metric for communication cost. Researchers have proposed alternative minimization objective functions that target per-processor communication metrics rather than global totals. An example is UMPa~\cite{umpa_2013}, which minimizes the maximum data sent by any single processor in a distributed memory setting. As Hendrickson noted, minimization objective functions that target communication costs are harder to optimize than the standard edge-cut, which partly explains their limited adoption.
\section{Thesis Goal}
 In this thesis we present a partition-aware communication volume objective function called \textit{NVOL} (see Chapters~\ref{cv_calculation}-\ref{chapter5}). NVOL minimizes the \textit{per-partition} maximum outgoing communication volume, aiming to reduce uneven communication loads. To be able to correctly calculate the partition-specific change in communication volume we also develop a novel algorithm in Chapter~\ref{cv_calculation} to perform this calculation. Using the METIS software as our baseline multilevel $k$-way partitioner, we present a new implementation of a partition-aware objective function and address two questions: 
\begin{itemize}
	\item \textbf{RQ1}: Can a partition-aware objective function produce more balanced communication loads than a global one?
	\item \textbf{RQ2}: Can the performance of a parallel application be improved in practice by using a partition-aware objective function?
\end{itemize}

NVOL is implemented in the refinement stage of the multilevel $k$-way algorithm. After an initial partition is performed, the partitions are gradually refined by moving vertices among partitions. NVOL gathers information about how every partition's communication volume is affected by a vertex move, allowing it to correctly evaluate how much partition $a$ will change if a vertex $v$ is moved to another partition $b$. Using this information, NVOL prioritizes moves that improve the per-partition maximum. No move that increases the maximum volume held by a single partition is accepted. This means that in some cases NVOL might even accept moves that increase the total communication volume, as long as they decrease the per-partition maximum. The goal is to provide more balanced partitions, avoiding overloading some for the sake of a better total communication volume.

\section{Thesis Structure}
\textbf{Chapter 2:} \textit{Background}
\begin{quote}
	Chapter 2 covers the fundamentals of graph partitioning. It surveys some of the algorithms in common use today, along with the main application areas in which they are applied.
\end{quote}
\textbf{Chapter 3:} \textit{Multilevel $k$-way Graph Partitioning}
\begin{quote}
	Chapter 3 presents the multilevel $k$-way graph partitioning algorithm on which NVOL is built. This approach offers a good balance among partition quality and runtime, making it a popular choice for large graphs. The chapter focuses on the standard formulation, which uses edge-cut; the following chapters then extend the algorithm to support volume-based objective functions such as NVOL.
\end{quote}
\textbf{Chapter 4:} \textit{Communication Volume Calculation}
\begin{quote}
	Chapter 4 explains how communication volume change is evaluated for volume-based objective functions. It covers the calculation of both total communication volume and per-partition volume, beginning with formal definitions of the calculations followed by worked examples. The final sections describe how these calculations are integrated into the multilevel $k$-way algorithm.
\end{quote}
\textbf{Chapter 5:} \textit{Greedy Refinement Using Communication Volume}
\begin{quote}
	Chapter 5 covers the refinement stage and explains how it is implemented for volume-based objective functions. Building on Chapter 4, it shows how partitions can be refined using communication volume objective functions. The total communication volume objective function and NVOL are each treated in their own section, and both conclude with a simplified algorithm representing the implementation.
\end{quote}
\textbf{Chapter 6:} \textit{METIS with NVOL and Supporting Tooling}
\begin{quote}
	Chapter 6 presents the tools used in this thesis to run and analyze the multilevel $k$-way partitioning algorithm with NVOL. It introduces METIS, the partitioning framework that NVOL extends, and describes its input and output formats. It also presents a small collection of Python scripts developed alongside this work for preprocessing input graphs and analyzing the per-partition communication volumes produced by the experiments.
\end{quote}
\textbf{Chapter 7:} \textit{Results}
\begin{quote}
	Chapter 7 presents the results of running NVOL across a variety of 	graphs and partition counts. It compares the partitions across multiple metrics and estimates of the resulting communication time among partitions produced by NVOL and METIS.
\end{quote}
\textbf{Chapter 8:} \textit{Conclusions and Future Work}
\begin{quote}
	Chapter 8 concludes the thesis and evaluates the contributions of this project. It revisits the initial goals set out in Chapter 1 and discusses how well they have been fulfilled. It discusses some of the current shortcomings and how they might be addressed, and introduces new ideas that would build naturally on the NVOL implementation.
\end{quote}